\documentclass[11pt]{article}
\newcommand{\ListaTitulo}{Lista 11 --- Quadrado Latino e Greco-Latino}
% Preâmbulo compartilhado — MATD48, Listas de Exercícios 2026
% Usado por Lista{NN}.tex e Gabarito{NN}.tex via % Preâmbulo compartilhado — MATD48, Listas de Exercícios 2026
% Usado por Lista{NN}.tex e Gabarito{NN}.tex via % Preâmbulo compartilhado — MATD48, Listas de Exercícios 2026
% Usado por Lista{NN}.tex e Gabarito{NN}.tex via \input{preamble}

\usepackage[utf8]{inputenc}
\usepackage[T1]{fontenc}
\usepackage[brazil]{babel}
\usepackage{fullpage}
\usepackage{amsmath,amssymb}
\usepackage{enumitem}
\usepackage{xcolor}
\usepackage{fancyhdr}
\usepackage{titlesec}
\usepackage{listings}
\usepackage{hyperref}
\usepackage{booktabs}
\usepackage{graphicx}

\definecolor{matd48azul}{HTML}{0B4F6C}
\definecolor{matd48verde}{HTML}{2E8B57}

\titleformat{\section}{\normalfont\Large\bfseries\color{matd48azul}}{\thesection}{1em}{}
\titleformat{\subsection}{\normalfont\large\bfseries\color{matd48azul}}{\thesubsection}{1em}{}

\providecommand{\ListaTitulo}{Lista de Exercícios}

\setlength{\headheight}{14pt}
\pagestyle{fancy}
\fancyhf{}
\lhead{\small MATD48 --- Planejamento de Experimentos A}
\rhead{\small \ListaTitulo}
\cfoot{\thepage}
\renewcommand{\headrulewidth}{0.4pt}

\lstset{
  basicstyle=\ttfamily\small,
  breaklines=true,
  frame=single,
  columns=fullflexible,
  backgroundcolor=\color{gray!8},
  commentstyle=\color{matd48verde},
  keywordstyle=\color{matd48azul}\bfseries,
  language=R,
  extendedchars=true,
  literate=%
    {á}{{\'a}}1 {à}{{\`a}}1 {â}{{\^a}}1 {ã}{{\~a}}1
    {é}{{\'e}}1 {ê}{{\^e}}1 {í}{{\'i}}1 {ó}{{\'o}}1
    {ô}{{\^o}}1 {õ}{{\~o}}1 {ú}{{\'u}}1 {ç}{{\c c}}1
    {Á}{{\'A}}1 {À}{{\`A}}1 {Â}{{\^A}}1 {Ã}{{\~A}}1
    {É}{{\'E}}1 {Ê}{{\^E}}1 {Í}{{\'I}}1 {Ó}{{\'O}}1
    {Ô}{{\^O}}1 {Õ}{{\~O}}1 {Ú}{{\'U}}1 {Ç}{{\c C}}1
}

\hypersetup{colorlinks=true, linkcolor=matd48azul, urlcolor=matd48azul, citecolor=matd48azul}

\newcommand{\cabecalho}[2]{%
  \begin{center}
    {\Large\bfseries\color{matd48azul} #1}\\[0.3em]
    {\large #2}\\[0.3em]
    \rule{\linewidth}{0.6pt}
  \end{center}
  \vspace{0.5em}
}

\newenvironment{questao}[1]{\vspace{1em}\noindent\textbf{Questão #1.}\hspace{0.4em}}{\par}
\newenvironment{solucao}{\vspace{0.3em}\noindent\textit{Solução.}\hspace{0.4em}\color{black!85}}{\par\vspace{0.5em}}


\usepackage[utf8]{inputenc}
\usepackage[T1]{fontenc}
\usepackage[brazil]{babel}
\usepackage{fullpage}
\usepackage{amsmath,amssymb}
\usepackage{enumitem}
\usepackage{xcolor}
\usepackage{fancyhdr}
\usepackage{titlesec}
\usepackage{listings}
\usepackage{hyperref}
\usepackage{booktabs}
\usepackage{graphicx}

\definecolor{matd48azul}{HTML}{0B4F6C}
\definecolor{matd48verde}{HTML}{2E8B57}

\titleformat{\section}{\normalfont\Large\bfseries\color{matd48azul}}{\thesection}{1em}{}
\titleformat{\subsection}{\normalfont\large\bfseries\color{matd48azul}}{\thesubsection}{1em}{}

\providecommand{\ListaTitulo}{Lista de Exercícios}

\setlength{\headheight}{14pt}
\pagestyle{fancy}
\fancyhf{}
\lhead{\small MATD48 --- Planejamento de Experimentos A}
\rhead{\small \ListaTitulo}
\cfoot{\thepage}
\renewcommand{\headrulewidth}{0.4pt}

\lstset{
  basicstyle=\ttfamily\small,
  breaklines=true,
  frame=single,
  columns=fullflexible,
  backgroundcolor=\color{gray!8},
  commentstyle=\color{matd48verde},
  keywordstyle=\color{matd48azul}\bfseries,
  language=R,
  extendedchars=true,
  literate=%
    {á}{{\'a}}1 {à}{{\`a}}1 {â}{{\^a}}1 {ã}{{\~a}}1
    {é}{{\'e}}1 {ê}{{\^e}}1 {í}{{\'i}}1 {ó}{{\'o}}1
    {ô}{{\^o}}1 {õ}{{\~o}}1 {ú}{{\'u}}1 {ç}{{\c c}}1
    {Á}{{\'A}}1 {À}{{\`A}}1 {Â}{{\^A}}1 {Ã}{{\~A}}1
    {É}{{\'E}}1 {Ê}{{\^E}}1 {Í}{{\'I}}1 {Ó}{{\'O}}1
    {Ô}{{\^O}}1 {Õ}{{\~O}}1 {Ú}{{\'U}}1 {Ç}{{\c C}}1
}

\hypersetup{colorlinks=true, linkcolor=matd48azul, urlcolor=matd48azul, citecolor=matd48azul}

\newcommand{\cabecalho}[2]{%
  \begin{center}
    {\Large\bfseries\color{matd48azul} #1}\\[0.3em]
    {\large #2}\\[0.3em]
    \rule{\linewidth}{0.6pt}
  \end{center}
  \vspace{0.5em}
}

\newenvironment{questao}[1]{\vspace{1em}\noindent\textbf{Questão #1.}\hspace{0.4em}}{\par}
\newenvironment{solucao}{\vspace{0.3em}\noindent\textit{Solução.}\hspace{0.4em}\color{black!85}}{\par\vspace{0.5em}}


\usepackage[utf8]{inputenc}
\usepackage[T1]{fontenc}
\usepackage[brazil]{babel}
\usepackage{fullpage}
\usepackage{amsmath,amssymb}
\usepackage{enumitem}
\usepackage{xcolor}
\usepackage{fancyhdr}
\usepackage{titlesec}
\usepackage{listings}
\usepackage{hyperref}
\usepackage{booktabs}
\usepackage{graphicx}

\definecolor{matd48azul}{HTML}{0B4F6C}
\definecolor{matd48verde}{HTML}{2E8B57}

\titleformat{\section}{\normalfont\Large\bfseries\color{matd48azul}}{\thesection}{1em}{}
\titleformat{\subsection}{\normalfont\large\bfseries\color{matd48azul}}{\thesubsection}{1em}{}

\providecommand{\ListaTitulo}{Lista de Exercícios}

\setlength{\headheight}{14pt}
\pagestyle{fancy}
\fancyhf{}
\lhead{\small MATD48 --- Planejamento de Experimentos A}
\rhead{\small \ListaTitulo}
\cfoot{\thepage}
\renewcommand{\headrulewidth}{0.4pt}

\lstset{
  basicstyle=\ttfamily\small,
  breaklines=true,
  frame=single,
  columns=fullflexible,
  backgroundcolor=\color{gray!8},
  commentstyle=\color{matd48verde},
  keywordstyle=\color{matd48azul}\bfseries,
  language=R,
  extendedchars=true,
  literate=%
    {á}{{\'a}}1 {à}{{\`a}}1 {â}{{\^a}}1 {ã}{{\~a}}1
    {é}{{\'e}}1 {ê}{{\^e}}1 {í}{{\'i}}1 {ó}{{\'o}}1
    {ô}{{\^o}}1 {õ}{{\~o}}1 {ú}{{\'u}}1 {ç}{{\c c}}1
    {Á}{{\'A}}1 {À}{{\`A}}1 {Â}{{\^A}}1 {Ã}{{\~A}}1
    {É}{{\'E}}1 {Ê}{{\^E}}1 {Í}{{\'I}}1 {Ó}{{\'O}}1
    {Ô}{{\^O}}1 {Õ}{{\~O}}1 {Ú}{{\'U}}1 {Ç}{{\c C}}1
}

\hypersetup{colorlinks=true, linkcolor=matd48azul, urlcolor=matd48azul, citecolor=matd48azul}

\newcommand{\cabecalho}[2]{%
  \begin{center}
    {\Large\bfseries\color{matd48azul} #1}\\[0.3em]
    {\large #2}\\[0.3em]
    \rule{\linewidth}{0.6pt}
  \end{center}
  \vspace{0.5em}
}

\newenvironment{questao}[1]{\vspace{1em}\noindent\textbf{Questão #1.}\hspace{0.4em}}{\par}
\newenvironment{solucao}{\vspace{0.3em}\noindent\textit{Solução.}\hspace{0.4em}\color{black!85}}{\par\vspace{0.5em}}


\begin{document}

\cabecalho{Lista de Exercícios 11}{Quadrado latino e quadrado greco-latino (Capítulo 4 / Aula 11)}

\noindent \textit{Entrega: até a próxima aula. Justifique todas as respostas — nestas listas, a
justificativa vale mais que a resposta final. O gabarito completo está em \texttt{Gabarito11.pdf}.}

\begin{questao}{1} \textbf{(Agricultura --- ANOVA de um quadrado latino)}

Um experimento com $n=5$ variedades de soja é conduzido em um quadrado latino $5\times5$, com
linha = faixa de declividade e coluna = faixa de umidade do talhão. O ajuste produziu
$SQ_{Total}=300$, $SQ_{Linha}=40$, $SQ_{Coluna}=35$ e $SQ_{Variedade}=150$.

\begin{enumerate}[label=(\alph*)]
  \item Complete a tabela de ANOVA (graus de liberdade de cada fonte, $SQ_{Erro}$, quadrados
    médios).
  \item Calcule $F_0$ para o teste do efeito de variedade.
  \item Sabendo que $F_{0{,}05;\,4;\,12}=3{,}26$, há evidência de diferença entre variedades?
\end{enumerate}
\end{questao}

\begin{questao}{2} \textbf{(Dedução --- graus de liberdade do erro em um quadrado latino replicado)}

Considere $p$ réplicas independentes de um quadrado latino de ordem $n$, cada réplica com suas
próprias linhas e colunas (não compartilhadas entre réplicas), mas com os mesmos $n$
tratamentos em todas as réplicas. O modelo tem termos para réplica, linha-dentro-de-réplica,
coluna-dentro-de-réplica e tratamento.

\begin{enumerate}[label=(\alph*)]
  \item Quantas observações há ao todo? Quantos graus de liberdade totais isso representa?
  \item Escreva os graus de liberdade consumidos por cada termo do modelo (réplica, linhas dentro
    de réplica, colunas dentro de réplica, tratamento).
  \item Subtraia a soma do item (b) do total do item (a) e mostre, com álgebra, que o resultado
    se simplifica para $gl_{Erro} = (n-1)\big[p(n-1)-1\big]$.
  \item Confira sua fórmula para $n=4$ e $p=2$: o valor deve coincidir com o que foi mostrado em
    aula.
\end{enumerate}
\end{questao}

\begin{questao}{3} \textbf{(Verificação --- validade de um quadrado latino)}

Um técnico propõe o seguinte arranjo $4\times4$ para 4 fertilizantes (A, B, C, D), com linhas =
talhões e colunas = datas de aplicação:
$$
\begin{array}{cccc}
A & B & C & D \\
B & A & D & C \\
C & D & A & B \\
D & C & A & B \\
\end{array}
$$

\begin{enumerate}[label=(\alph*)]
  \item Verifique, linha a linha e coluna a coluna, se este arranjo satisfaz a definição de
    quadrado latino. Se não satisfizer, aponte exatamente onde a definição é violada.
  \item Caso tenha encontrado um problema no item (a), proponha uma correção mínima (alterando o
    menor número possível de células) que restaure a validade do quadrado.
\end{enumerate}
\end{questao}

\begin{questao}{4} \textbf{(Quadrado greco-latino --- existência e graus de liberdade)}

\begin{enumerate}[label=(\alph*)]
  \item Calcule os graus de liberdade do erro, $(n-1)(n-3)$, para quadrados greco-latinos de
    ordem $n=5$ e $n=7$.
  \item Por que não existe um quadrado greco-latino de ordem $n=6$? A que resultado clássico esse
    fato se refere?
  \item Um pesquisador com $n=4$ fontes de bloqueio a controlar (por exemplo, linha, coluna e
    \emph{dois} fatores de tratamento adicionais) percebe que precisaria de um "quadrado
    hiper-greco-latino" com 3 conjuntos de letras sobrepostos, não apenas 2. Quantos graus de
    liberdade sobrariam para o erro nesse caso, com $n=4$? Esse desenho é praticamente viável?
    Justifique.
\end{enumerate}
\end{questao}

\begin{questao}{5} \textbf{(Psicologia --- quadrado latino para contrabalancear ordem)}

Um estudo de carga cognitiva expõe cada participante a $4$ tipos de tarefa (A, B, C, D) em $4$
sessões consecutivas. O pesquisador teme dois efeitos de confusão: diferenças individuais entre
participantes e efeito de \emph{prática/fadiga} (a ordem em que uma tarefa é apresentada — 1ª,
2ª, 3ª ou 4ª posição — pode, por si só, afetar o desempenho).

\begin{enumerate}[label=(\alph*)]
  \item Proponha um desenho em quadrado latino $4\times4$ para este estudo, identificando
    explicitamente o que fará o papel de linha, de coluna e de tratamento.
  \item Explique por que esse desenho, e não um DBCA com bloco = participante apenas, é apropriado
    aqui — que fonte de variação um DBCA simples deixaria de controlar?
  \item O desenho do item (a) assume ausência de interação entre ordem de apresentação e tipo de
    tarefa. Dê um exemplo concreto de como essa suposição poderia falhar neste estudo.
\end{enumerate}
\end{questao}

\begin{questao}{6} \textbf{(Aleatorização de Yates e análise de poder)}

\begin{enumerate}[label=(\alph*)]
  \item Um técnico decide sempre usar o mesmo quadrado latino $4\times4$ fixo (com o tratamento A
    sempre na diagonal principal) em todo experimento de sua empresa que peça esse desenho.
    Explique por que essa prática pode comprometer a validade da inferência, e descreva os quatro
    passos do procedimento de aleatorização de Yates (1933) que evitam o problema.
  \item Um pesquisador com $n=4$ tratamentos planeja usar quadrados latinos replicados e deseja
    poder $\ge0{,}80$ para detectar um efeito grande ($f=0{,}4$ de Cohen) a $\alpha=0{,}05$. A
    tabela abaixo (vista em aula) mostra o poder em função do número de réplicas $p$:
    $$
    \begin{array}{c|ccccc}
    p & 1 & 2 & 3 & 4 & 5 \\\hline
    \text{Poder} & 0{,}148 & 0{,}356 & 0{,}557 & 0{,}716 & 0{,}828
    \end{array}
    $$
    Quantas réplicas são necessárias? Que trade-off prático essa escolha envolve, em termos de
    unidades experimentais totais gastas?
\end{enumerate}
\end{questao}

\end{document}
